\subsection*{2.1 Equivalent Expressions}


\noindent\markforreview
\vspace{1em}

The equation \((ax + 3)(5x^2 - bx + 4) = 20x^3 - 9x^2 - 2x + 12\) is true for all \(x\), where \(a\) and \(b\) are constants. What is the value of \(ab\)?

\vspace{0.75em}

\choicebox{A}{18}
\choicebox{B}{20}
\choicebox{C}{24}
\choicebox{D}{40}


\vspace{2em}

\noindent\markforreview
\vspace{1em}

Which of the following expressions is(are) a factor of \(3x^2 + 20x - 63\)?

I. \(x - 9\)  

II. \(3x - 7\)

\vspace{0.75em}

\choicebox{A}{I only}
\choicebox{B}{II only}
\choicebox{C}{I and II}
\choicebox{D}{Neither I nor II}


\vspace{2em}


\noindent\markforreview
\vspace{1em}

If \(\frac{\sqrt{x^5} }{\sqrt[3]{x^4}}  = x^{\frac{a}{b}}\) for all positive values of \(x\), what is the value of \(\frac{a}{b}\)?



\newpage
\noindent\markforreview
\vspace{1em}

The equation \(\frac{2}{x-2} + \frac{3}{x+5} = \frac{rx + t}{(x-2)(x+5)}\) is true for all \(x > 2\), where \(r\) and \(t\) are positive constants. What is the value of \(rt\)?

\vspace{0.75em}

\choicebox{A}{-20}
\choicebox{B}{15}
\choicebox{C}{20}
\choicebox{D}{60}

\vspace{2em}
\noindent\markforreview
\vspace{1em}

For what value of \(x\) is the given expression \(\sqrt[5]{70n} \left(\sqrt[6]{70n}\right)^2\) equivalent to \((70n)^{30x}\), where \(n > 1\)?


\vspace{2em}

\noindent\markforreview
\vspace{1em}

The expression \((3x - 23)(19x + 6)\) is equivalent to the expression \(ax^2 + bx + c\), where \(a, b, c\) are constants. What is the value of \(b\)?





\newpage 
\noindent\markforreview
\vspace{1em}

Which expression is equivalent to \(\frac{4}{4x - 5} - \frac{1}{x + 1}?\)

\vspace{0.75em}

\choicebox{A}{\(\frac{1}{(x+1)(4x-5)}\)}
\choicebox{B}{\(\frac{3}{3x - 6}\)}
\choicebox{C}{\(-\frac{1}{(x+1)(4x - 5)}\)}
\choicebox{D}{\(\frac{9}{(x+1)(4x - 5)}\)}


\vspace{2em}

\noindent\markforreview
\vspace{1em}

In the expression \(\frac{x^2 - c}{x - b}\), \(b\) and \(c\) are positive integers.  
If the expression is equivalent to \(x + b\) and \(x \ne b\), which of the following could be the value of \(c\)?

\vspace{0.75em}

\choicebox{A}{4}
\choicebox{B}{6}
\choicebox{C}{8}
\choicebox{D}{10}
\newpage
\noindent\markforreview
\vspace{1em}

The expression \(0.36x^2 + 0.63x + 1.17\) can be rewritten as  
\(a(4x^2 + 7x + 13)\), where \(a\) is a constant.  
What is the value of \(a\)?



\vspace{2em}


\noindent\markforreview
\vspace{1em}

In the expression \(3(2x^2 + px + 8) - 16x(p + 4)\),  
\(p\) is a constant. This expression is equivalent to  
\(6x^2 - 155x + 24\). What is the value of \(p\)?

\vspace{0.75em}

\choicebox{A}{-3}
\choicebox{B}{7}
\choicebox{C}{13}
\choicebox{D}{155}


\vspace{2em}


\noindent\markforreview
\vspace{1em}

The expression \(\frac{x^{-2}y^{\frac{1}{2}}}{x^{\frac{1}{3}}y^{-1}}\), where \(x > 1\) and \(y > 1\), is equivalent to which of the following?

\vspace{0.75em}

\choicebox{A}{\(\frac{\sqrt{y}}{\sqrt[3]{x^2}}\)}
\choicebox{B}{\(\frac{y\sqrt{y}}{\sqrt[3]{x^{2}}}\)}
\choicebox{C}{\(\frac{y\sqrt{y}}{x\sqrt{x}}\)}
\choicebox{D}{\(\frac{y\sqrt{y}}{x^{2}\sqrt[3]{x}}\)}


\newpage

\noindent\markforreview
\vspace{1em}

The expression \(4x^2 + bx - 45\), where \(b\) is a constant,  
can be rewritten as \((hx + k)(x + j)\), where \(h, k,\) and \(j\) are integers.  
Which of the following must be an integer?

\vspace{0.75em}

\choicebox{A}{\(\frac{b}{h}\)}
\choicebox{B}{\(\frac{b}{k}\)}
\choicebox{C}{\(\frac{45}{h}\)}
\choicebox{D}{\(\frac{45}{k}\)}

\vspace{2em}


\noindent\markforreview
\vspace{1em}

Let \(f(x) = x^3 - 9x\), \(g(x) = x^2 - 2x - 3\).  
Which of the following expressions is \(\frac{f(x)}{g(x)}\) equivalent to, for \(x > 3\)?

\vspace{0.75em}

\choicebox{A}{\(\frac{1}{x + 1}\)}
\choicebox{B}{\(\frac{x + 3}{x + 1}\)}
\choicebox{C}{\(\frac{x(x - 3)}{x + 1}\)}
\choicebox{D}{\(\frac{x(x + 3)}{x + 1}\)}

\newpage
\noindent\markforreview
\vspace{1em}

The expression \(\frac{1}{3}x^2 - 2\) can be rewritten as \(\frac{1}{3}(x - k)(x + k)\),  
where \(k\) is a positive constant. What is the value of \(k\)?

\vspace{0.75em}

\choicebox{A}{2}
\choicebox{B}{6}
\choicebox{C}{\(\sqrt{2}\)}
\choicebox{D}{\(\sqrt{6}\)}


\vspace{2em}
\noindent\markforreview
\vspace{1em}

Which of the following is equivalent to \(\left(a+\frac{b}{2}\right)^2\) ?

\vspace{0.75em}

\choicebox{A}{\(a^2 + \frac{b^2}{2}{}\)}
\choicebox{B}{\(\frac{a^2 + b^2}{4}\)}
\choicebox{C}{\(\frac{a^2 + ab}{2} + \frac{b^2}{2}\)}
\choicebox{D}{\(a^2 + ab +\frac{ b^2}{4}\)}

\newpage

\noindent\markforreview
\vspace{1em}

Which expression is equivalent to \(\frac{y + 12}{x - 8} + \frac{y(x - 8)}{x^2 y - 8xy}\) ?



\vspace{0.75em}

\choicebox{A}{\(\frac{xy + y + 4}{x^2 y - 16x + 64xy}\)}
\choicebox{B}{\(\frac{xy + 9 + 12}{x^2 y - 8xy + x - 8}\)}
\choicebox{C}{\(\frac{xy^2 + 13xy - 8y}{x^2 y - 8xy}\)}
\choicebox{D}{\(\frac{xy^2 + 13xy - 8y}{x^2 y - 16x^2 y + 64xy}\)}


\vspace{2em}


\noindent\markforreview
\vspace{1em}

One of the factors of \(2x^3 + 42x^2 + 208x\) is \(x + b\),  
where \(b\) is a positive constant. What is the smallest possible value of \(b\)?

\vspace{0.75em}



\vspace{2em}

\noindent\markforreview
\vspace{1em}

The equation \(\frac{x^2 + 6x - 7}{x + 7} = ax + d\) is true for all \(x \ne -7\),  
where \(a\) and \(d\) are integers. What is the value of \(a + d\)?

\vspace{0.75em}

\choicebox{A}{-6}
\choicebox{B}{-1}
\choicebox{C}{0}
\choicebox{D}{1}


\newpage

\noindent\markforreview
\vspace{1em}

Which of the following expressions is equivalent to \(\frac{x^2 - 2x - 5}{x - 3}\) ?

\vspace{0.75em}

\choicebox{A}{\(x - 5-\frac{  20}{x - 3}\)}
\choicebox{B}{\(x - 5 -\frac{ 10}{x - 3}\)}
\choicebox{C}{\(x + 1 -\frac{ 8}{x - 3}\)}
\choicebox{D}{\(x + 1 - \frac{2}{x - 3}\)}

\vspace{2em}

\noindent\markforreview
\vspace{1em}

The expression \((7532 + 100y^2) + 10(10y^2 - 110)\)  
can be written in the form \(ay^2 + b\), where \(a\) and \(b\) are constants.  
What is the value of \(a + b\)?

\vspace{0.75em}


\newpage


\noindent\markforreview
\vspace{1em}

The expression \(6 \sqrt[5]{3^5x^{45}} \cdot \sqrt[8]{2^8x}\) is equivalent to \(ax^b\),  
where \(a\) and \(b\) are positive constants and \(x > 1\).  
What is the value of \(a + b\)?

\vspace{0.75em}



\vspace{2em}

\noindent\markforreview
\vspace{1em}

If \(a = c + d\), which of the following is equivalent to the expression  
\(x^2 - c^2 - 2cd - d^2\)?

\vspace{0.75em}

\choicebox{A}{\((x + a)^2\)}
\choicebox{B}{\((x - a)^2\)}
\choicebox{C}{\((x + a)(x - a)\)}
\choicebox{D}{\(x^2 - ax - a^2\)}


\vspace{2em}

\noindent\markforreview
\vspace{1em}

If \(a\) and \(c\) are positive numbers, which of the following  
is equivalent to \(\sqrt{(a + c)^3} \cdot \sqrt{a + c}\) ?

\vspace{0.75em}

\choicebox{A}{\(a + c\)}
\choicebox{B}{\(a^2 + c^2\)}
\choicebox{C}{\(a^2 + 2ac + c^2\)}
\choicebox{D}{\(a^2 c^2\)}



\newpage
